\documentclass[a4paper,12pt]{article}
\usepackage{graphicx, setspace, array, xcolor, mathtools, contour, tikz, amsthm, setspace, amsmath,amsfonts,amssymb, subcaption, fancyhdr, lipsum,multicol, geometry, gensymb}
\usepackage{colortbl}
\geometry{a4paper, margin=0.75in}
\contourlength{0.5pt}
\usetikzlibrary{patterns}
\usepackage[english]{babel}
\renewcommand{\qedsymbol}{$\blacksquare$}
\definecolor{darkgreen}{rgb}{0.0, 0.5, 0.0}
\pagestyle{fancy}
\fancyhf{}
\renewcommand{\headrulewidth}{0.4pt}
\fancyhead[L]{\rightmark}
\fancyhead[R]{\thepage}
\title{Modern Algebra}
\author{Assignment 13\\ \\ Yousef A. Abood\\ \\ ID: 900248250}
\date{December 2025}
\setlength{\parindent}{0pt}
\singlespacing
\parskip=1mm
\setcounter{section}{12}
\setcounter{subsection}{1}
\onehalfspacing
\begin{document}
\maketitle
\noindent\makebox[\linewidth]{\rule{15cm}{0.4pt}}
\section{Integral Domains}
\subsection*{Problem 3}
%  
%  Show that a commutative ring with the cancellation property ­(under
% multiplication) has no zero-divisors.
\begin{proof}
    Suppose we have a commutative ring $R$ with the cancellation property, that is,\\ if $ba=bc$ then $a=c$ where $a,b,c \in R$. Suppose that $a,b$ are non-zero elements. For the sake of contradiction, we assume that $R$ has a zero-divisor $b$ such that $ba=0.$ 
    We know that $b\cdot a=b\cdot 0$. By the cancellation property, we cancel $b$, so $a = 0$, which contradicts our assumption that $a \ne 0$. Hence, a commutative ring cannot have a zero-divisor.
\end{proof}
\subsection*{Problem 4}
%   List all zero-divisors in Z20. Can you see a relationship between the
% zero-divisors of Z20 and the units of Z20?
The set of zero-divisors of $\mathbb{Z}_{20}$, call it $S$, is: $\{2,4,5,6,8,10,12,14,15,18,16\}$.\\
We know that the set of units $U(\mathbb{Z}_{20})=\{1,3,7,9,11,13,17,19\}$
\\ Thus, $S \cap U(\mathbb{Z}_{20})=\emptyset.$ The relationship between the two sets is that they are disjoint.
\subsection*{Problem 9}
% Find elements a, b, and c in the ring Z x Z x Z such that ab, ac, and
% bc are zero-divisors but abc is not a zero-divisor.
Consider $a=(2,0,1), b=(3,1,0), c=(0,4,1)$ and observe that $ab=(6,0,0), bc=(0,4,0), ac=(0,0,1)$. Clearly, $ab,bc,ac$ are non-zero elements. In otherwords, none on them is the additive idintity of the ring $\mathbb{Z} \oplus \mathbb{Z}\oplus \mathbb{Z}.$ Observe that 
\begin{align*}
    ab \cdot bc &= (6 \cdot 0, 0 \cdot 4, 0 \cdot 0)=(0,0,0)\\
    bc\cdot ac &= (0 \cdot 0, 4 \cdot 0, 0 \cdot 1)=(0,0,0)\\
    a \cdot b \cdot c &=(2 \cdot 3 \cdot 0, 0 \cdot 1 \cdot 4, 1 \cdot 0 \cdot 1)=(0,0,0).
\end{align*}
Thus, $abc=(0,0,0)$ which is the additive identity, so it is not a zero-divisor. Hence, we found elements $a,b,c \in \mathbb{Z} \oplus \mathbb{Z}\oplus \mathbb{Z}$ such that $ab,bc,ac$ are zero divisors, and $abc$ is not a zero divisor.
\subsection*{Problem 18}
%  A ring element a is called an idempotent if a2 5 a. Prove that the
% only idempotents in an integral domain are 0 and 1.
\begin{proof}
    Suppose we have an integral domain $D$ and $a$ is an element in $D.$ If $a$ is idempotent, then $a^2=a.$ Observe that
    \begin{align*}
        a^2   &=a\\
        a^2-a &=a-a\\
        a^2-a &=0_D\\
        a(a-1)&=0_D
    \end{align*}
    Then either $a-1=0$ or $a=0$, so $a=1$ or $a=0.$ Hence, if there are idempotents in an integral domain then they are only $1$ or $0$.
\end{proof}
\subsection*{Problem 32}
%  Let R 5 {0, 2, 4, 6, 8} under addition and multiplication modulo
% 10. Prove that R is a field.
Consider the ring $\mathbb{Z}_{10}.$ Clearly, $R \subset Z_{10}$. We need to show it is a subring. First, we show it is closed under subtraction, observe the following Cayley table:
\begin{table}[h]
    \centering
    \begin{tabular}{|c|c|c|c|c|c|}
        \hline
      $-$ & \textbf{0} & \textbf{2} & \textbf{4} & \textbf{6} & \textbf{8} \\
      \hline
      \textbf{0} & 0 & 8 & 6 & 4 & 2 \\
      \hline
      \textbf{2} & 2 & 0 & 8 & 6 & 4 \\
      \hline
      \textbf{4} & 4 & 2 & 0 & 8 & 6 \\
      \hline
      \textbf{6} & 6 & 4 & 2 & 0 & 8 \\
      \hline
      \textbf{8} & 8 & 6 & 4 & 2 & 0 \\
      \hline
    \end{tabular}
\end{table}
\\Next, we show it is closed under multiplication, observe the following Cayley table:
\begin{table}[h]
    \centering
    \begin{tabular}{|c|c|c|c|c|c|}
        \hline
      $\cdot$ & \textbf{0} & \textbf{2} & \textbf{4} & \textbf{6} & \textbf{8} \\
      \hline
      \textbf{0} & 0 & 0 & 0 & 0 & 0 \\
      \hline
      \textbf{2} & 0 & 4 & 8 & 2 & 6 \\
      \hline
      \textbf{4} & 0 & 8 & 6 & 4 & 2 \\
      \hline
      \textbf{6} & 0 & 2 & 4 & 6 & 8 \\
      \hline
      \textbf{8} & 0 & 6 & 2 & 8 & 4 \\
      \hline
    \end{tabular}
\end{table}
Hence, $R$ is a ring. From the previous table, we can see that if $a \in R$, then $a \cdot 6 =a$, so $6$ is the mulitplicative identity. We see that $6$ is in every row except the zero row, thus, every element has a mulitplicative inverse. For commutativity, suppose that we have an index for each entry in the table $ij$, where $1\le i,j \le 5$. From the multiplication table, the entries in the index $ij$ are the same in the entries $ji.$ Therefore, $R$ is a field.

\subsection*{Problem 57}
\begin{itemize}
    \item [a)] We factor the expression to get $(x-2)(x-3)=0$. We know that $\mathbb{Z}_7$ is an integral domain since $7$ is prime number. Then, $x$ is either $2$ or $3$. So we have two solution for the equation $x^2 - 5x + 6=0$ in $\mathbb{Z}_7$.
    \item [b)] We factor the expression to get $(x-2)(x-3)=0$. We know that $2, 3 \in \mathbb{Z}_8$, so $x=2$ or $x=3$ are solutions for the equation. $\mathbb{Z}_8$ is not an integral domain, then we can have other solutions for this equation other than $x=2,x=3.$ However, there is no element $x \in \mathbb{Z}_8$ such that $(x\ne 2 \wedge x\ne 3) \wedge ((x-2) \cdot (x-3)=0).$ Thus, the equation $x^2 - 5x + 6=0$ has only two solutions in $\mathbb{Z}_8$
\end{itemize}
\end{document}